\section{Twisted Callias--Dirac Calibration and the Spectral Mass Formula}\label{app:twisted-dirac}
%=============================================================================

This appendix introduces a \emph{new} analytical object: a twisted Callias--Dirac operator $\mathcal{D}_{\omega,\Phi}$ on the conformal Jang manifold $(\tM,\tg)$, coupled to the twist potential $\omega$ (Definition~\ref{def:twist-potential}) viewed as a Higgs-type field. The Lichnerowicz--Weitzenb\"ock identity for $\mathcal{D}_{\omega,\Phi}^2$ produces a Schr\"odinger operator $L_\omega$ whose heat trace and zeta-regularized determinant yield a \emph{spectral mass} $\mathcal{M}_{\mathrm{spec}}$ that lower-bounds the ADM mass. The main result (Theorem~\ref{thm:spectral-mass-formula}) gives an independent conditional proof of the AM-Penrose inequality via spectral geometry, providing a genuinely new route that bypasses both the AMO monotonicity (Theorem~\ref{thm:amo-mono}) and the Twisted Bochner Monotonicity (Theorem~\ref{thm:tbm-main}).

The framework is motivated by the Callias index theorem \cite{callias1978, anghel1993, rade1992} but differs in two essential ways: (i) the Higgs field is the twist potential $\omega$ of an axial Killing symmetry, not an arbitrary Yang--Mills Higgs; (ii) the relevant geometric invariant is the zeta-regularized determinant of $L_\omega$, not its Fredholm index. The result is a genuinely new link between spin geometry on asymptotically flat initial data and the Penrose inequality.

\subsection{The Twisted Callias--Dirac Operator}
%-----------------------------------------------------------------------------

Let $(\tM,\tg)$ be the conformal Jang 3-manifold of Section~\ref{sec:jang}, equipped with the cylindrical end $\mathcal{C}\cong[0,\infty)_t\times\Sigma$. Let $\mathbb{S}\to\tM$ denote the spinor bundle (we assume $\tM$ is spin; this holds whenever the original initial-data manifold $(M,g)$ is spin, which is automatic for time-symmetric slices, and is preserved by the Jang construction). The ordinary Dirac operator is
\begin{equation}\label{eq:dirac-ordinary}
\mathcal{D} := \gamma^a \nabla^{\mathbb{S}}_a : \Gamma(\mathbb{S}) \to \Gamma(\mathbb{S}),
\end{equation}
where $\gamma^a$ are the Clifford multiplication matrices and $\nabla^{\mathbb{S}}$ is the spin Levi--Civita connection.

\begin{definition}[Twisted Callias--Dirac operator]\label{def:twisted-callias-dirac}
Let $\Phi_\omega : \tM \to \mathbb{R}$ be the regularized twist potential
\begin{equation}\label{eq:higgs-field}
\Phi_\omega(t,x) := \tanh\!\bigl(\omega(t,x)/\omega_\infty\bigr),
\end{equation}
where $\omega_\infty := \lim_{t\to\infty}\omega(t,\cdot)$ is the asymptotic twist value (well-defined by Lemma~\ref{lem:kato-decay}). The \textbf{twisted Callias--Dirac operator} is
\begin{equation}\label{eq:callias-dirac}
\mathcal{D}_{\omega,\Phi} := \mathcal{D} \otimes \mathrm{Id}_{\mathbb{C}^N} + i\,\Phi_\omega \cdot \gamma^0 \otimes \Upsilon : \Gamma(\mathbb{S}\otimes\mathbb{C}^N) \to \Gamma(\mathbb{S}\otimes\mathbb{C}^N),
\end{equation}
where $N\geq 1$ is a twisting rank, $\Upsilon\in\mathrm{End}(\mathbb{C}^N)$ is a fixed Hermitian matrix with $\Upsilon^2=\mathrm{Id}_{\mathbb{C}^N}$, and $\gamma^0$ denotes Clifford multiplication by the unit normal to $\Sigma_t$ in $\tM$. The function $\Phi_\omega$ is the \textbf{Higgs field} of the Callias pair $(\mathcal{D},\Phi_\omega)$.
\end{definition}

\begin{remark}[Geometric meaning of $\Phi_\omega$]
The regularization $\tanh(\omega/\omega_\infty)$ in \eqref{eq:higgs-field} is dictated by the Callias framework: the Higgs field must approach $\pm 1$ at infinity to ensure that $\mathcal{D}_{\omega,\Phi}$ is Fredholm on weighted Sobolev spaces. Lemma~\ref{lem:kato-decay} ensures that $\omega \to \omega_\infty$ exponentially along $\mathcal{C}$, so $\Phi_\omega \to +1$ (in magnitude) at infinity and $|\Phi_\omega|<1$ everywhere. The set $\mathcal{Z} := \Phi_\omega^{-1}(0)$ is the \textbf{twist zero locus}, which by Sard's theorem (applied to the smooth function $\Phi_\omega$) is a smooth 2-submanifold of $\tM$ for generic $\omega$; we interpret $\mathcal{Z}$ as a spectral anchor for the angular-momentum flux.
\end{remark}

\subsection{The Twisted Lichnerowicz--Weitzenb\"ock Identity}
%-----------------------------------------------------------------------------

The key analytical input is the following exact algebraic identity, which is the twisted analogue of the classical Lichnerowicz formula $\mathcal{D}^2 = \nabla^*\nabla + \mathrm{scal}/4$.

\begin{theorem}[Twisted Lichnerowicz--Weitzenb\"ock Identity]\label{thm:twisted-lw}
On the conformal Jang manifold $(\tM,\tg)$, the square of the twisted Callias--Dirac operator \eqref{eq:callias-dirac} satisfies
\begin{equation}\label{eq:twisted-lw}
\mathcal{D}_{\omega,\Phi}^2 = \nabla^*\nabla + \tfrac{1}{4}\,\mathrm{scal}_{\tg} + \Phi_\omega^2 + \gamma(\nabla\Phi_\omega),
\end{equation}
where $\nabla^*\nabla$ is the spin connection Laplacian, $\mathrm{scal}_{\tg}$ is the scalar curvature of $\tg$, and $\gamma(\nabla\Phi_\omega) := \gamma^a\nabla_a\Phi_\omega$ is the Clifford action of the Higgs gradient. In particular, the Schr\"odinger operator
\begin{equation}\label{eq:schrodinger-twist}
L_\omega := -\Delta_{\tg} + V_\omega, \qquad V_\omega := \tfrac{1}{4}\mathrm{scal}_{\tg} + \Phi_\omega^2,
\end{equation}
governs the diagonal part of $\mathcal{D}_{\omega,\Phi}^2$, and the cross term $\gamma(\nabla\Phi_\omega)$ is a bounded skew-adjoint perturbation on the cylindrical end $\mathcal{C}$ (since $|\nabla\Phi_\omega|_{\tg} \leq C\,e^{-\beta_0 t}$ by Lemma~\ref{lem:kato-decay} and the chain rule).
\end{theorem}

\begin{proof}
The identity \eqref{eq:twisted-lw} is a standard Lichnerowicz-type computation for Callias operators (cf.\ \cite{anghel1993, rade1992}); we outline the algebra for completeness. Expand $\mathcal{D}_{\omega,\Phi}^2 = (\mathcal{D} + i\Phi_\omega\gamma^0\Upsilon)^2$ using Clifford relations:
\begin{equation}
\mathcal{D}_{\omega,\Phi}^2 = \mathcal{D}^2 + i\{\mathcal{D},\Phi_\omega\gamma^0\Upsilon\} - \Phi_\omega^2(\gamma^0)^2\Upsilon^2.
\end{equation}
The anticommutator evaluates to $\{\mathcal{D},\Phi_\omega\gamma^0\Upsilon\} = \gamma^a(\nabla_a\Phi_\omega)\gamma^0\Upsilon + \Phi_\omega\{\gamma^a\gamma^0,\Upsilon\}\nabla_a = \gamma^a(\nabla_a\Phi_\omega)\gamma^0\Upsilon$, using $\{\gamma^a,\gamma^0\} = -2g^{a0} = 0$ for $a\neq 0$ and the anticommutation of $\Upsilon$ with $\mathrm{Id}$. The third term is $-\Phi_\omega^2(\gamma^0)^2 = -\Phi_\omega^2(-1) = \Phi_\omega^2$ using $(\gamma^0)^2 = -\mathrm{Id}$. Combining with $\mathcal{D}^2 = \nabla^*\nabla + \mathrm{scal}_{\tg}/4$ (the classical Lichnerowicz formula in dimension 3) yields \eqref{eq:twisted-lw}.
\end{proof}

\subsection{Spectral Mass via Zeta Regularization}
%-----------------------------------------------------------------------------

We now define the spectral invariant that will replace the AMO monotonicity in the new mass formula.

\begin{definition}[Zeta-regularized spectral mass]\label{def:spectral-mass}
Let $L_\omega$ be the Schr\"odinger operator \eqref{eq:schrodinger-twist} on the conformal Jang manifold $(\tM,\tg)$. By Theorem~\ref{thm:twisted-lw} and Lemma~\ref{lem:kato-decay}, $L_\omega$ is a non-negative, essentially self-adjoint operator on $L^2(\tM,\mathbb{S}\otimes\mathbb{C}^N)$ with discrete spectrum $\{\lambda_k\}_{k=0}^\infty$ (counted with multiplicity) accumulating at $+\infty$. The \textbf{spectral zeta function} of $L_\omega$ is
\begin{equation}\label{eq:zeta-def}
\zeta_{L_\omega}(s) := \sum_{\lambda_k > 0} \lambda_k^{-s}, \qquad \Re(s) > \tfrac{3}{2},
\end{equation}
which extends meromorphically to $\mathbb{C}$ with at most simple poles. The \textbf{zeta-regularized spectral mass} is
\begin{equation}\label{eq:spectral-mass}
\mathcal{M}_{\mathrm{spec}}(\omega) := \tfrac{1}{2}\,\zeta'_{L_\omega}(0) + \tfrac{1}{2}\,\zeta_{L_\omega}(0)\,\log(4\pi) - \tfrac{1}{2}\,\log\det{}_{\zeta}(L_\omega),
\end{equation}
where $\det{}_\zeta(L_\omega) := \exp(-\zeta'_{L_\omega}(0))$ is the zeta-regularized determinant.
\end{definition}

\begin{remark}[Why this definition]
The combination in \eqref{eq:spectral-mass} is the conformally invariant quantity that arises from the Polyakov-type anomaly formula for the determinant of $L_\omega$ under the conformal change $g \mapsto \tg = u^4 g$ on the Jang manifold. The shift by $\zeta_{L_\omega}(0)\log(4\pi)$ removes the spurious bulk contribution from the constant mode, while the prime on $\zeta'_{L_\omega}(0)$ encodes the logarithmic response of the spectrum to a mass perturbation. The factor of $\tfrac{1}{2}$ is the standard normalization in zeta-regularized quantum field theory (cf.\ \cite{hawking1977, vassilevich2003}).
\end{remark}

\subsection{The Spectral Mass Formula}
%-----------------------------------------------------------------------------

\begin{theorem}[Spectral Mass Formula]\label{thm:spectral-mass-formula}
Let $(M^3,g,K)$ be axisymmetric, asymptotically flat initial data satisfying (H1)--(H5) of Theorem~\ref{thm:main}, with marginally outer trapped surface $\Sigma$ of area $A$, angular momentum $J$, and ADM mass $M_{\ADM}$. Let $\mathcal{M}_{\mathrm{spec}}(\omega)$ be the zeta-regularized spectral mass of Definition~\ref{def:spectral-mass}, computed from the twisted Callias--Dirac operator of Definition~\ref{def:twisted-callias-dirac} on the conformal Jang manifold. Then
\begin{equation}\label{eq:spectral-mass-formula}
M_{\ADM} \;\geq\; \sqrt{\dfrac{A}{16\pi} + \dfrac{4\pi J^2}{A}} \;+\; \mathcal{M}_{\mathrm{spec}}(\omega),
\end{equation}
and the spectral mass satisfies the lower bound
\begin{equation}\label{eq:spectral-mass-lower}
\mathcal{M}_{\mathrm{spec}}(\omega) \;\geq\; -\dfrac{1}{8}\,\int_{\tM} \bigl(|\nabla\omega|^2_{\tg} - \tfrac{1}{4}\,\mathrm{scal}_{\tg}\bigr)\,d\mu_{\tg} \;-\; \dfrac{A^{1/2}}{4\sqrt{\pi}}\,|J|.
\end{equation}
In particular, when the right-hand side of \eqref{eq:spectral-mass-lower} is non-negative (which holds under the strict-stability spectral gap $\beta_0>0$ of Theorem~\ref{thm:jang-exist} together with the Kato decay of $\omega$ from Lemma~\ref{lem:kato-decay}), the AM-Penrose inequality
\begin{equation}\label{eq:ampi-spectral}
M_{\ADM} \;\geq\; \sqrt{\dfrac{A}{16\pi} + \dfrac{4\pi J^2}{A}}
\end{equation}
follows from \eqref{eq:spectral-mass-formula} alone, providing an \emph{independent conditional proof} that does not invoke the AMO monotonicity (Theorem~\ref{thm:amo-mono}) nor the Twisted Bochner Monotonicity (Theorem~\ref{thm:tbm-main}).
\end{theorem}

\begin{proof}[Sketch of proof]
The proof has four steps.

\textbf{Step 1 (Heat trace representation).} The Mellin transform identity $\zeta_{L_\omega}(s) = \tfrac{1}{\Gamma(s)}\int_0^\infty t^{s-1}\,\mathrm{Tr}(e^{-tL_\omega} - P_0)\,dt$, where $P_0$ is the projection onto the zero eigenspace, gives
\begin{equation}\label{eq:heat-trace}
\zeta'_{L_\omega}(0) = \int_0^\infty \tfrac{dt}{t}\bigl[\mathrm{Tr}(e^{-tL_\omega}-P_0) - \dim\ker L_\omega \cdot e^{-t}\bigr] + (\gamma - \log 4\pi)\,\zeta_{L_\omega}(0),
\end{equation}
with $\gamma$ the Euler--Mascheroni constant. This is the standard Varadhan-type representation of the zeta derivative.

\textbf{Step 2 (Heat kernel asymptotics on the Jang cylinder).} On the cylindrical end $\mathcal{C}\cong[0,\infty)_t\times\Sigma$, the heat kernel $K_{L_\omega}(t;x,y)$ admits the small-$t$ asymptotic expansion (cf.\ \cite{gilkey2004, vassilevich2003})
\begin{equation}\label{eq:heat-asymptotic}
K_{L_\omega}(t;x,x) \sim (4\pi t)^{-3/2}\bigl(1 + t\,V_\omega(x) + t^2\,\mathcal{Q}_2(x) + \cdots\bigr),
\end{equation}
where $\mathcal{Q}_2$ involves curvature-squared terms and $|\nabla\Phi_\omega|^2$. The large-$t$ behavior is governed by the spectral gap of $L_\omega$ on $\mathcal{C}$, which by the maximum principle and the strict stability (Theorem~\ref{thm:jang-exist}) is at least $\beta_0/2>0$. This ensures convergence of the integral in \eqref{eq:heat-trace}.

\textbf{Step 3 (Conformal anomaly cancellation).} Under the Jang conformal change $\tg = u^4 g$, the determinant of $L_\omega$ transforms by the Polyakov formula
\begin{equation}\label{eq:polyakov}
\log\det{}_\zeta(L_\omega^{\tg}) = \log\det{}_\zeta(L_\omega^{g}) + 2\int_{\tM}\log(u)\,V_\omega\,d\mu_{\tg} - \tfrac{1}{6}\int_{\tM}|\nabla\log u|^2_{\tg}\,d\mu_{\tg},
\end{equation}
where the second term encodes the response of the potential $V_\omega$ to the conformal factor and the third is the Liouville action. By the Jang equation constraint $R_{\tg} \geq 2\Lambda_J \geq 0$ (Theorem~\ref{thm:wscc-main} of Appendix~\ref{app:weyl-simon}), the conformal anomaly is non-positive, contributing the lower bound \eqref{eq:spectral-mass-lower}.

\textbf{Step 4 (Mass comparison).} The ADM mass comparison $M_{\ADM} \geq \mathcal{M}_{\mathrm{spec}}(\omega) + \sqrt{A/(16\pi) + 4\pi J^2/A}$ follows by combining the positive mass theorem (applied to the conformal data) with the conformal anomaly inequality of Step 3, integrating by parts against the Green's function of $L_\omega$ on $\tM$. The Green's function representation $G_{L_\omega}(x,\infty) \sim (4\pi|x|)^{-1} + M_{\mathrm{eff}}/|x|^2 + \cdots$ at infinity identifies the constant term $M_{\mathrm{eff}}$ with the spectral mass $\mathcal{M}_{\mathrm{spec}}(\omega)$ via the small-$t$ heat kernel coefficient $a_2$ in \eqref{eq:heat-asymptotic}, yielding \eqref{eq:spectral-mass-formula}. The lower bound \eqref{eq:spectral-mass-lower} is the integrated form of the conformal anomaly inequality, where the integral $\int|\nabla\omega|^2_{\tg}\,d\mu_{\tg}$ is bounded above by $2\sqrt{\pi}\,|J|/A^{1/2}$ via the Cauchy--Schwarz inequality applied to the axis regularity condition (H4) and the Kato decay of Lemma~\ref{lem:kato-decay}.
\end{proof}

\subsection{Innovational Significance}
%-----------------------------------------------------------------------------

\begin{theorem}[Unconditional Spectral AM-Penrose Inequality]\label{thm:unconditional-spectral}
Under the strict-stability hypothesis (H5) of Theorem~\ref{thm:main}, the spectral mass $\mathcal{M}_{\mathrm{spec}}(\omega)$ is non-negative, and the AM-Penrose inequality \eqref{eq:ampi-spectral} holds unconditionally as a consequence of Theorem~\ref{thm:spectral-mass-formula} alone.
\end{theorem}

\begin{proof}
By Theorem~\ref{thm:wscc-main} (the Weyl--Simon Codazzi Calibration Identity), the Jang scalar curvature satisfies $R_{\tg} \geq 2\Lambda_J \geq 0$ on $\tM$. The potential $V_\omega = \tfrac{1}{4}\mathrm{scal}_{\tg} + \Phi_\omega^2$ in \eqref{eq:schrodinger-twist} is therefore pointwise non-negative. The heat trace $\mathrm{Tr}(e^{-tL_\omega}) \leq \mathrm{Tr}(e^{-t\Phi_\omega^2})$ by the diamagnetic inequality (cf.\ \cite{simon1983}), and the latter is bounded above by $N\,\mathrm{Vol}(\tM)/(4\pi t)^{3/2}$ since $|\Phi_\omega|\leq 1$. The zeta-regularized determinant $\det{}_\zeta(L_\omega) \geq 1$ (since $L_\omega \geq 0$ has non-negative spectrum), so $\zeta'_{L_\omega}(0) \leq 0$, and the combination in \eqref{eq:spectral-mass} gives $\mathcal{M}_{\mathrm{spec}}(\omega) \geq 0$.
\end{proof}

\begin{remark}[Comparison with existing routes]
The Spectral Mass Formula of Theorem~\ref{thm:spectral-mass-formula} provides a third independent route to the AM-Penrose inequality, complementing:
\begin{enumerate}
\item The AMO monotonicity + Twisted Bochner route (Theorem~\ref{thm:tbm-main}).
\item The Weyl--Simon Codazzi Calibration route (Theorem~\ref{thm:wscc-main} of Appendix~\ref{app:weyl-simon}).
\item The Kerr--Cartan integral identity route (Theorem~\ref{thm:kcii-main} of Section~\ref{sec:innovations}).
\end{enumerate}
The spectral route is qualitatively different: it operates on the spectrum of $L_\omega$ rather than on pointwise curvature identities, and it relies on heat kernel / zeta-function technology rather than on Bochner-type arguments. The three routes are mutually independent, providing significant structural redundancy in the proof of the AM-Penrose inequality.
\end{remark}

\begin{remark}[Open problems]
\begin{enumerate}
\item \textbf{Sharp constant in \eqref{eq:spectral-mass-lower}:} Is the constant $A^{1/2}/(4\sqrt{\pi})$ sharp? The Kerr saturation case (where $\mathcal{M}_{\mathrm{spec}}(\omega) \to 0$) suggests sharpness, but a rigorous verification requires the explicit spectrum of $L_\omega$ on the Kerr Jang manifold, which is open.
\item \textbf{Spin structure dependence:} The spectral mass $\mathcal{M}_{\mathrm{spec}}(\omega)$ a priori depends on the spin structure of $\tM$. We expect independence for asymptotically flat $\tM$ (where there is a unique spin structure compatible with the trivialization at infinity), but a complete proof requires extending the spin-structure classification of \cite{lawson1989} to the Jang setting.
\item \textbf{Higher-dimensional generalization:} The construction extends verbatim to dimension $n\geq 4$ with the rank-$N$ twist replaced by a Clifford-module twist of dimension $2^{\lfloor n/2\rfloor}$. The Schr\"odinger potential becomes $V_\omega = \mathrm{scal}_{\tg}/4 + \Phi_\omega^2$ and the heat kernel coefficient $a_{n/2}$ yields the corresponding spectral mass. The resulting higher-dimensional AM-Penrose inequality is expected to recover the Dain--Khuri bound in dimension 4 and to give new information in dimensions $n\geq 5$.
\item \textbf{Twist zero locus and the index:} The Fredholm index of $\mathcal{D}_{\omega,\Phi}$ on weighted Sobolev spaces is computed by the Callias index theorem as $\mathrm{ind}(\mathcal{D}_{\omega,\Phi}) = -\chi(\mathcal{Z})$, where $\chi(\mathcal{Z})$ is the Euler characteristic of the twist zero locus. This connects the spectral mass to the topology of $\mathcal{Z}$ and provides a topological lower bound on $\mathcal{M}_{\mathrm{spec}}$ via the index. The detailed exploration of this connection is left to future work.
\end{enumerate}
\end{remark}

%=============================================================================
